\section{A conductor-seventeen indicator balance}
\label{sec:q17-indicator-balance}

\begin{theorem}\label{thm:q17-five-exclusions}
The five rings
\[
R8\_2\_33,\quad R8\_2\_36,\quad R8\_2\_76,\quad
R8\_2\_412,\quad R8\_2\_446
\]
have no complex fusion categorification.
\end{theorem}

The same short indicator identity excludes all five rings, including the
noncommutative cases. No enumeration of possible center rows or
associator equations is needed.

\begin{lemma}\label{lem:q17-conductor-equality}
Let $S,T$ be the normalized unitary modular representation of a
pseudo-unitary Drinfeld center, and let $P$ project onto the $T$-fixed
space. Suppose that a $T$-fixed vector $w$ has diagonal-torus character
the primitive quadratic character modulo $17$. If a known subspace
$L$ of the $T$-fixed space contains $w$ and
\[
\operatorname{proj}_L Sw=-\frac{w}{\sqrt{17}},
\]
then $PSw=-w/\sqrt{17}$ and the cyclic representation generated by $w$
factors through level $17$.
\end{lemma}
\begin{proof}
The compression bound in \cite[Lemma~4.1]{SchDimensionBlock} is already
attained on $L$, so the full projection cannot have any additional
coordinate. For completeness, the equality case also gives the asserted
level bound. By the congruence theorem \cite{NgCongruence}, decompose the
finite representation into local irreducible tensor factors. At $17$, an
irreducible representation with an upper-unipotent-fixed vector embeds
in the primitive-vector permutation representation over
$\mathbb Z/17^a\mathbb Z$. In the quadratic torus eigenspace, the
maximal singular subspace of the compression is spanned by the functions
\[
F(p,q)=\chi_{17}(q)[q\text{ is a unit}],\qquad
G(p,q)=\chi_{17}(-p)[q=0\pmod {17}].
\]
This is the equality subspace in the proof of the cited local bound.
Both functions are pullbacks from level $17$. An irreducible constituent
containing an equality vector therefore factors through that level.
At every other prime the character is trivial; equality in the bound
one makes the vector fixed by both upper and lower unipotents, and
hence by the entire local group. In a tensor product of contractions,
the maximal singular subspace is the tensor product of the local maximal
singular subspaces, as follows by diagonalizing the positive operators
$A^*A$. Thus all constituents supporting $w$ factor through level $17$.
The full equality argument is also recorded in the evidence note
\texttt{compression/CONDUCTOR\_EQUALITY\_LEMMA.md}.
\end{proof}

\begin{lemma}[Saturated Galois indicator balance]
\label{lem:saturated-galois-indicator}
Let $\mathcal C$ be pseudo-unitary, with simple dimensions in a real
quadratic field $K$ of conductor $f$ and irrational global dimension
$D$. Suppose every center simple $A$ outside $I(1)$ with
$\dim(A)/D\notin\mathbb Q$ satisfies $\theta_A^f=1$.
Let $\ell:K\longrightarrow\mathbb Q$ be a nonzero rational linear map
with $\ell(D)=0$. If $k$ is coprime to $f$ and
$\zeta_f\mapsto\zeta_f^k$ fixes $K$, then
\begin{equation}\label{eq:saturated-galois-indicator}
\sum_{i\ne0}\ell(\dim X_i)\nu_k(X_i)=0.
\end{equation}
\end{lemma}
\begin{proof}
Choose a common cyclotomic twist field $\mathbb Q(\zeta_N)$, with
$f\mid N$. The Chinese remainder theorem gives an exponent congruent
to $k$ modulo $f$ and coprime to $N$: at each prime dividing $N$ but
not $f$, impose the additional congruence $1$. The corresponding
Galois automorphism $\tau$ fixes $K$ and raises $f$th roots to their
$k$th powers.

For $z_A=\dim A$, induction reciprocity and the first trace identity
\cite[Theorem~2.5]{OstrikPiv3} give
\[
\sum_A\ell(z_A)z_A\theta_A=D\ell(1).
\]
All coefficients belong to $K$. The terms with $z_A/D\in\mathbb Q$
vanish since $\ell(z_A)=0$. In every other term outside $I(1)$ the
twist is an $f$th root by hypothesis; inside $I(1)$ it is one.
Applying $\tau$ therefore replaces every remaining $\theta_A$ by
$\theta_A^k$, while fixing the coefficients and right side. The general
induction indicator formula \cite[Theorem~4.1]{NgIndicators} now gives
\[
D\sum_i\ell(\dim X_i)\nu_k(X_i)=D\ell(1).
\]
Cancel $D$ and the unit term. The argument proves an exact complex
identity, without taking real parts or relative traces, and does not
assume that the source ring is commutative.
\end{proof}

\begin{proof}[Proof of Theorem~\ref{thm:q17-five-exclusions}]
Every source ring has four invertibles, forming a group $G$ of order
four, and four noninvertibles $\mathcal O$, all of dimension
\[
d=4+\sqrt{17},\qquad d^2=1+8d,\qquad
D=4+4d^2=8+32d=136+32\sqrt{17}.
\]
By \cite[Lemma~3.1]{SchNear}, any complex categorification is Galois
conjugate to a pseudo-unitary one. This two-FP-dimension theorem does
not assume commutativity. Work with the conjugate and its canonical
positive spherical structure.

The source tensors give
\[
\Omega=\sum_X XX^*=8\,1+8\sum_{\rho\in\mathcal O}\rho,
\qquad
\det(xI-L_\Omega)=(x-8)^6(x^2-272x+1088).
\]
Here $L_\Omega$ is the Casimir multiplication matrix; in the
noncommutative cases we do not identify it with the full induction
Gram matrix. If an irreducible ring representation $E$ has dimension
$e_E$, the Casimir eigenvalue is $e_E f_E$ with multiplicity $e_E^2$
\cite[Remark~2.11]{OstrikPiv3}. Thus $D$ and
$D'=136-32\sqrt{17}$ are the codegrees of one-dimensional
representations, and every other ordinary codegree $f_E=8/e_E$ is
rational.

By \cite[Theorem~2.13]{OstrikPiv3}, the corresponding center constituent
$A_E$ of $I(1)$ has multiplicity $e_E$ and dimension $D/f_E$.
Let $J$ correspond to $D'$. Then
\[
[I(1):J]=1,\qquad \dim J=D/D'=d^2.
\]
All these constituents have twist one. Also $F(I(1))=\Omega$ implies
that none contains a nonunit invertible in its restriction.

Normalize $S_{00}=1/D$. The multiplicity vector $m$ of $I(1)$ satisfies
$Sm=m$ \cite[Proposition~2.8]{OstrikPiv3}. Therefore its support is
Galois-stable with positive Galois signs. A Galois automorphism fixing
$K=\mathbb Q(\sqrt{17})$ sends the central unit to the unique constituent
of $I(1)$ of dimension one; one flipping $K$ sends it to $J$, the unique
constituent of dimension $D/D'$. Thus
$w=e_1-e_J$ has quadratic torus character $\chi_{17}$.

The known twist-one block sends $w$ to
\[
\left(\frac1D-\frac1{D'}\right)w=-\frac{w}{\sqrt{17}}.
\]
Indeed its $1,J$ block has diagonal $1/D$ and off-diagonal $1/D'$,
and the other entries in these two columns agree because their
codegrees are rational. This remains true when $e_E=2$.
Lemma~\ref{lem:q17-conductor-equality} now gives
$PSw=-w/\sqrt{17}$ and level $17$ on the cyclic representation.
For a center simple $A$ of dimension $z$, the coordinate of $Sw$ is
\[
\frac zD-\sigma\!\left(\frac zD\right).
\]
If $z/D$ is nonrational and $A$ is outside $I(1)$, its twist is therefore
a primitive seventeenth root: it has order dividing $17$, and twist
one is excluded by the exact formula for $PSw$. Apply Lemma~\ref{lem:saturated-galois-indicator} with
\[
\ell(a+b\sqrt{17})=4a-17b,\qquad k=2.
\]
Indeed $\ell(D)=0$, and $2$ is a quadratic residue modulo $17$, so
$\zeta_{17}\mapsto\zeta_{17}^2$ fixes $K$. The source weights are
$\ell(1)=4$ on the invertibles and $\ell(d)=-1$ on the
noninvertibles. We obtain the common balance
\begin{equation}\label{eq:q17-indicator-balance}
\boxed{\displaystyle
\sum_{\rho\in\mathcal O}\nu_2(\rho)
 =4\sum_{g\in G\setminus\{1\}}\nu_2(g).}
\end{equation}
No restrictions on twists of constituents of $I(g)$, no relative
moment hull, and no enumeration of center rows are needed.

A group of order four has an odd number of nonunit involutions: one
for $C_4$, three for $C_2^2$. Its nontrivial pointed indicator sum is
therefore a nonzero odd integer, so the right side of
\eqref{eq:q17-indicator-balance} has absolute value at least four.
The left side has absolute value at most the number of self-dual
noninvertibles. The exact source values are
\[
\begin{array}{c|c|c|c}
\text{ring}&G&\text{self-dual noninvertible labels}&\text{number}\\\hline
R8\_2\_33&C_4&5,6&2\\
R8\_2\_36&C_4&\varnothing&0\\
R8\_2\_76&C_4&\varnothing&0\\
R8\_2\_412&C_2^2&6,7&2\\
R8\_2\_446&C_2^2&\varnothing&0
\end{array}
\]
Each row contradicts the balance, proving all five exclusions over
$\mathbb C$.
\end{proof}

The source-bound verifier \path{verify_q17_indicator_balance.py}
checks the original tensors, hashes, associativity, FP dimensions,
pointed groups, duality lists, Casimir polynomial, and the displayed
arithmetic. It uses only exact integer operations and runs identically
under standard and optimized Python. The proof and independent audits
are stored in \path{evidence/unified_20261016/rank8/}, including
\texttt{Q17\_INDICATOR\_BALANCE\_PROOF.md}.
